% Paper 2 — Devanagari handwriting generation with latent diffusion.
% LNCS format (ICDAR/ICFHR target). Compile with XeLaTeX/LuaLaTeX for
% Devanagari examples; see paper1/main.tex for the fontspec pattern.
\documentclass{llncs}

\usepackage{graphicx}
\usepackage{booktabs}
\usepackage{amsmath}
\usepackage{url}
\usepackage{iftex}
\ifPDFTeX
  \newcommand{\dn}[1]{\texttt{[Devanagari]}}
\else
  \usepackage{fontspec}
  \newfontfamily\devafont[Script=Devanagari]{Noto Sans Devanagari}
  \newcommand{\dn}[1]{{\devafont #1}}
\fi

\begin{document}

\title{Glyph-Conditioned Latent Diffusion for Devanagari Handwritten
Word Generation}
\titlerunning{Devanagari Handwriting Generation with Latent Diffusion}

\author{Samir Wagle\inst{1} \and Manish~...\inst{1}}
\institute{Tribhuvan University, Kathmandu, Nepal\\
\email{samir@redtab.xyz}}

\maketitle

\begin{abstract}
Diffusion-based handwriting generation has advanced rapidly for Latin
scripts, but no latent-diffusion generator exists for Devanagari — an
abugida whose conjunct ligatures, dependent vowel signs, and connecting
headline make verbatim text rendering substantially harder than for
alphabetic scripts. We present the first latent diffusion system for
Devanagari handwritten word generation. Text content is enforced through a
ControlNet branch conditioned on font-rendered glyph images, sidestepping
the failure of CLIP-based text encoders on non-Latin orthography. We
evaluate with generation fidelity (FID), content fidelity (CER of an
independent Devanagari HTR model on generated images), and a conjunct-aware
out-of-vocabulary protocol measuring generalization to unseen consonant
clusters. Finally, we show downstream utility: mixing generated words into
HTR training data improves recognition on IIIT-INDIC-HW-WORDS Hindi. We
release the model, training code, and evaluation protocol, and outline an
akshara-tokenized conditioning scheme as the path to writer-styled
generation.
\keywords{Handwriting generation \and Latent diffusion \and ControlNet
\and Devanagari \and Data augmentation.}
\end{abstract}

\section{Introduction}

Generating realistic handwritten text images has two practical payoffs:
training-data augmentation for handwritten text recognition (HTR) in
low-resource scripts, and personalized text rendering. Latent diffusion
models (LDMs) now dominate this task for Latin
scripts~\cite{nikolaidou2023wordstylist,nikolaidou2024diffusionpen,dai2024onedm},
yet a survey of the literature shows \emph{no} diffusion-based handwriting
generator for Devanagari or any Indic script; prior Devanagari generation
work is limited to GANs over isolated characters~\cite{nepaligan2022}.

The gap is not incidental. Latin-script diffusion generators condition on
text via character embeddings or CLIP-style encoders. For Devanagari both
break down: CLIP tokenizers have effectively no Devanagari coverage, and
character-sequence conditioning must handle conjunct formation — the visual
form of \dn{क्} + \dn{ष} is the fused ligature \dn{क्ष}, not a
concatenation of glyphs. A generator that spells rather than ligates
produces text no literate reader accepts.

Our contributions:
\begin{enumerate}
  \item \textbf{The first Devanagari handwriting LDM.} A ControlNet branch
  over Stable Diffusion, conditioned on font-rendered glyph images, which
  delegates orthography (conjunct formation, matra placement, shirorekha)
  to a Unicode shaping engine and lets diffusion model stroke style and
  texture.
  \item \textbf{A script-aware evaluation protocol}: FID for visual
  fidelity; CER/AER of an independent TrOCR-based Devanagari recognizer on
  generated images for content fidelity; and an out-of-vocabulary split
  stratified by conjunct content, testing generalization to unseen
  clusters.
  \item \textbf{Downstream utility}: dose-response of mixing generated
  words into HTR training (companion study), the first augmentation result
  for Indic HTR from a diffusion generator.
\end{enumerate}

\section{Related Work}

\paragraph{Diffusion handwriting generation.}
WordStylist~\cite{nikolaidou2023wordstylist} trains a compact LDM from
scratch with character-embedding cross-attention and writer embeddings on
IAM. DiffusionPen~\cite{nikolaidou2024diffusionpen} adds a few-shot style
encoder; One-DM~\cite{dai2024onedm} achieves one-shot style imitation;
Diffusion Brush extends generation to full text lines. All target Latin
scripts; none address abugida orthography.

\paragraph{Glyph conditioning for visual text.}
GlyphControl~\cite{yang2023glyphcontrol} and successors condition scene-text
generation on rendered glyph images, showing that explicit visual glyph
guidance outperforms text-encoder conditioning for accurate spelling. We
bring this insight to handwriting: for a script where a shaping engine
already solves ligation, rendered-glyph conditioning transfers orthographic
correctness for free.

\paragraph{Devanagari generation.}
Existing work generates isolated Nepali characters with
GANs~\cite{nepaligan2022} at low resolution without text conditioning —
class-conditional synthesis, not verbatim word generation. To our knowledge
no prior system generates arbitrary Devanagari words in handwritten style.

\section{Method}

\subsection{Architecture}

We build on Stable Diffusion v1.5's VAE latent space and UNet, freezing
both, and train a ControlNet~\cite{zhang2023controlnet} branch initialized
from the UNet encoder. The conditioning image renders the target word in
Noto Sans Devanagari — white on black, foreground-cropped, centered at
$256{\times}256$ — so the Unicode shaping engine (HarfBuzz) resolves
conjuncts, matra positioning, and the shirorekha before diffusion begins.
The prompt pathway receives a fixed template with a phonetic
transliteration of the word; because CLIP cannot represent Devanagari, the
prompt contributes style context (``handwritten, paper, ink'') while the
ControlNet branch alone carries orthographic content. Training: 28k steps,
LR $1.5{\times}10^{-5}$, fp16, on IIIT-INDIC-HW-WORDS Hindi word images
resized to $256{\times}256$.

\subsection{Why Not Character Embeddings (Yet)}

WordStylist-style character-embedding conditioning must \emph{learn}
conjunct formation from data. In IIIT-INDIC-HW-WORDS Hindi, conjunct
frequency follows a long tail; rare clusters may appear only a handful of
times — a data regime where learned ligation fails silently. Rendered-glyph
conditioning guarantees orthographic correctness for every representable
cluster, at the cost of style controllability. Section~\ref{sec:future}
outlines the akshara-tokenized hybrid we see as the next step.

\section{Evaluation Protocol}

\paragraph{Visual fidelity.} FID between generated and real word images
(test-split references), computed at generation resolution.
\paragraph{Content fidelity.} An independently trained Devanagari TrOCR
recognizer (from our companion study) reads each generated image; we report
CER and akshara error rate against the intended text. This measures whether
generations are \emph{readable as the right word}, not merely plausible.
\paragraph{Conjunct-stratified OOV.} We generate (i) in-vocabulary words,
(ii) OOV words composed of seen aksharas, and (iii) OOV words containing
conjuncts unseen in training, and compare content fidelity across strata.
\paragraph{Downstream augmentation.} Generated words mixed into HTR
training at $+10/25/50\%$ of real data; recognition CER on the official
test split.

\section{Results}

% TODO(experiments pending — run after Paper 1 sweep completes):
% 1. Generate eval sets with paper1/experiments/generate_synthetic.py
%    (IV / OOV-seen-akshara / OOV-unseen-conjunct vocab files).
% 2. FID: clean-fid against real test images.
% 3. Content fidelity: evaluate.py --data-dir over generated sets with the
%    best Paper-1 adapter as the judge.
% 4. Augmentation: run_augmentation.sh dose-response table.
% 5. Qualitative grid: real vs generated for 8 words incl. heavy conjuncts.
\emph{[Experiments in progress; tables and qualitative figures to be
inserted from the evaluation pipeline.]}

\section{Limitations and Future Work}
\label{sec:future}

Rendered-glyph conditioning trades style diversity for orthographic
correctness: generations track the conditioning layout, and writer identity
is not controllable. The path forward is akshara-level tokenized
conditioning — embeddings over orthographic syllables rather than
codepoints, giving the model ligation-aware content control — combined with
a style encoder over writer exemplars in the DiffusionPen mold. The public
Hindi word benchmark exposes no writer identities, so style supervision
requires either the original dataset release's writer metadata or
unsupervised style pseudo-labels; we leave this to future work.

\section{Conclusion}

We presented the first latent-diffusion handwriting generator for
Devanagari, a script-aware evaluation protocol measuring what Latin-script
metrics miss (conjunct generalization, content readability), and evidence
that diffusion-generated data improves Devanagari HTR. The gap between
Latin and Indic handwriting generation is now open to systematic study.

\bibliographystyle{splncs04}
\bibliography{references}

\end{document}
